\scanpage{23}
\mspage{22}
\begin{center}- 53 -\end{center}
holomorphe dans $P$ (comme fonction réciproque de $e^z$), on a $\int_\gamma\frac{dz}{z}=\log(\gamma(b))-\log(\gamma(a))$, et par suite $\unsure{\Im}(\int_\gamma\frac{dz}{z})=\mathrm{Am}(\gamma(b))-\mathrm{Am}(\gamma(a))$, en désignant par $\mathrm{Am}(z)$ la mesure de l'amplitude de $z$ appartenant à l'intervalle $\left]-\pi,+\pi\right[$. Par abus de langage, on dit encore que $\unsure{\Im}(\int_\gamma\frac{dz}{z})$ pour une route \uline{quelconque} $\gamma$ contenue dans $\unsure{\complement\{0\}|\mathbb C\setminus\{0\}}$, est la \uline{variation de l'amplitude de} $z$ \uline{le long de} $\gamma$. Si $\gamma$ est un circuit d'indice $n$ par rapport à $0$, la ``variation de l'amplitude de $z$ '' le long de $\gamma$ est donc $2n\pi$.

Les notions précédentes relatives aux chemins \uline{carrossables} peuvent s'étendre aux chemins \uline{quelconques}. Soit $\gamma$ un chemin achevé quelconque défini dans un intervalle $I=[\alpha,\beta]$ à valeurs dans $\mathbb C$, et soit $a$ un point de $\mathbb C$ n'appartenant pas à $\gamma(I)$. Il existe une déformation de $\gamma$ en un chemin carrossable $\gamma'$ \uline{dans} $\complement\{a\}$ (§ 2, prop.1) conservant les extrémités de $\gamma$ ; en outre, si $\gamma'$ et $\gamma''$ sont deux chemins carrossables déformés de $\gamma$ dans $\complement\{a\}$ par de telles déformations, $\gamma'$ et $\gamma''$ peuvent être déformés l'un en l'autre dans $\complement\{a\}$ par une déformation de même type ; en raison du cor. du th.1 du § 2, la ``variation de l'amplitude'' de $z-a$ le long de $\gamma'$ et le long de $\gamma''$ sont égales. Cette variation ne dépend donc que du chemin $\gamma$ initial, et on l'appelle encore la \uline{variation de l'amplitude de} $z-a$ \uline{le long de} $\gamma$.

Il résulte aussitôt de ces définitions que les variations de l'amplitude de $z-a$ le long de $\gamma$ et le long de l'opposé de $\gamma$ sont \uline{opposées} ; si $\gamma_1$ et $\gamma_2$ sont deux chemins dans $\complement\{a\}$ tels que la juxtaposition de ces chemins soit définie, la variation de l'amplitude de $z-a$ le long de $\unsure{\gamma_1\vee\gamma_2|\gamma_1\wedge\gamma_2}$ est la \uline{somme} des variations de l'amplitude de $z-a$ le long de $\gamma_1$ et le long de $\gamma_2$.
